\section{问题分析}

\subsection{问题一的分析}

问题一的本质是允许硬模块旋转 $90^\circ$ 的自由轮廓二维矩形装箱问题，需要同时确定模块方向、平面坐标与外接轮廓。三组实例含 $100$--$300$ 个模块，排列、旋转和坐标的完全枚举不可行，因此将自由轮廓搜索转化为候选轮廓宽度搜索，并分别用 Skyline 与 MaxRects 构造布局。题目规定先比较轮廓面积、仅在面积相同时再比较长短边比，故采用词典序目标保持这一判优顺序。

\subsection{问题二的分析}

相较问题一，问题二同时加入由给定死区率确定的固定正方形轮廓、模块间网络连接和外部 Terminal。模块不仅需要满足旋转、边界与两两非重叠约束，其中心位置还通过网络共同决定总半周长线长（HPWL），因而本问成为几何可行性与互连线长相互耦合的组合优化问题。

若直接联合搜索模块位置、旋转状态、非重叠关系与 HPWL，既难从离散可行域中获得稳定的优化方向，也难在搜索过程中持续维持几何可行。为此，本文沿“网络驱动目标位置—几何合法化—局部线长优化”的主线求解：先以二次线长松弛提取网络意义上的目标中心，再将目标中心转化为满足固定轮廓与非重叠约束的合法布局，最后对合法化引入的线长损失进行局部修正。该分层处理使网络信息能够进入几何放置过程，同时保留对布局可行性的显式控制。

\subsection{问题三的分析}

问题三的建模逻辑分为三层。第一，全部硬模块的总面积固定，而正方形死区率仅随轮廓边长 $S$ 单调变化，因此最小化死区率等价于最小化 $S$。这使面积压缩目标能够转化为边长明确、步长自然的整数搜索问题。

第二，题目要求先压缩轮廓、再控制互连线长，两项目标具有明确的先后次序。本文据此采用 $\operatorname{lexmin}(S,\mathcal{L})$，先比较轮廓边长，再在相同边长下比较总 HPWL，避免人为设置面积与线长的加权系数。

第三，从模块总面积和最大模块边长给出的理论下界开始，逐整数搜索正方形轮廓，并以问题二的已知可行边界为上界；在当前搜索策略首次成功构造全部模块布局的边长下，复用问题二的网络目标中心、合法化与局部重插入方法优化 HPWL。由此形成“理论下界---整数轮廓搜索---紧轮廓线长优化”的求解链条，并直接比较面积利用率改善与互连代价。

\subsection{问题四的分析}

问题四的建模路线可压缩为四步：
\begin{enumerate}
  \item \textbf{识别外接矩形失效。}L 型、T 型模块含有凹槽，若仍用矩形外接框判断碰撞，会把凹槽中的空闲区域误判为已占用，从而排除本应可行的紧凑布局；
  \item \textbf{建立真实形状集合。}将模块表示为实际占用区域的集合，生成四向旋转状态，并以旋转、平移后的集合是否相交判断重叠；
  \item \textbf{进行完整有限枚举。}本问仅有 $4$ 个整数尺寸模块，搜索空间有限，故枚举旋转状态与放置位置，以回溯方法寻找可行布局，无需采用随机启发式方法；
  \item \textbf{用面积下界判定最优。}模块总面积是任意外接轮廓面积的理论下界；若具体可行构造达到该下界，即可严格证明其为全局最优。
\end{enumerate}
